% 20260909 Science Quiz mathematics Problem set.

% Verified against OpenAI's official research announcement dated September 8, 2026.
\begin{problem}{1}
2026년 9월 8일 OpenAI는 AI가 생성한 writeup과 Lean formal proof를 공개하며 한 Clay Millennium Prize Problem에 대한 solution을 proposed했다고 발표했다. 이 문제의 이름을 쓰시오.
\end{problem}

\begin{solution}
OpenAI가 2026년 9월 8일 solution을 proposed했다고 발표한 Clay Millennium Prize Problem은
\[
    \boxed{\text{Navier--Stokes existence and smoothness problem}}
\]
이다. 발표에는 AI-generated writeup과 Lean formal proof가 함께 공개되었다.
\end{solution}

\begin{problem}{2}
수열 $(a_n)$이
\[
    a_0=2\cos\left(\frac{2\pi}{7}\right),
    \qquad
    a_{n+1}=a_n^2-2
\]
을 만족한다. $a_{2026}$을 구하시오.
\end{problem}

\begin{solution}
배각공식
\[
    (2\cos\theta)^2-2=2\cos(2\theta)
\]
에 의해 귀납적으로
\[
    a_n
    =
    2\cos\left(\frac{2^{n+1}\pi}{7}\right)
\]
을 얻는다. $2$의 거듭제곱을 $14$로 나눈 나머지는 주기 $3$을 가지며,
\[
    2^{2027}\equiv4\pmod{14}.
\]
따라서
\[
    \boxed{a_{2026}=2\cos\left(\frac{4\pi}{7}\right)}.
\]
\end{solution}

\begin{problem}{3}
\[
    \mathbf{A}\in\R^{2\times3},
    \qquad
    \mathbf{B}\in\R^{3\times2}
\]
이고
\[
    \mathbf{A}\mathbf{B}
    =
    \begin{pmatrix}
        2&1\\
        1&2
    \end{pmatrix}
\]
라고 하자. $3\times3$ 행렬 $\mathbf{B}\mathbf{A}$의 characteristic polynomial
\[
    p_{\mathbf{B}\mathbf{A}}(\lambda)
\]
를 구하시오.
\end{problem}

\begin{solution}
$\mathbf{A}\mathbf{B}$와 $\mathbf{B}\mathbf{A}$는 $0$이 아닌 eigenvalue들을 algebraic multiplicity까지 동일하게 갖는다. 먼저
\[
    \begin{aligned}
        p_{\mathbf{A}\mathbf{B}}(\lambda)
        &=
        \det\!\begin{pmatrix}
            \lambda-2&-1\\
            -1&\lambda-2
        \end{pmatrix}\\
        &=(\lambda-2)^2-1\\
        &=(\lambda-1)(\lambda-3).
    \end{aligned}
\]
따라서 $3\times3$ 행렬 $\mathbf{B}\mathbf{A}$의 eigenvalue는 $0,1,3$이고,
\[
    \boxed{
        p_{\mathbf{B}\mathbf{A}}(\lambda)
        =
        \lambda(\lambda-1)(\lambda-3)
    }.
\]
\end{solution}

% Verified against Bryce Putman's August 2026 counterexample and Jorik Jooken's human-checkable proof.
\begin{problem}{4}
2026년 8월 Bryce Putman이 $112$개의 꼭짓점을 가진 connected simple bridgeless cubic graph의 명시적 counterexample을 공개했고, Jorik Jooken이 이어서 human-checkable proof를 제시했다. 이 결과로 반박된 graph-theory conjecture의 통용 명칭을 쓰시오.
\end{problem}

\begin{solution}
반박된 conjecture는
\[
    \boxed{\text{Petersen Coloring Conjecture}}
\]
이다. 2026년 8월 Putman이 $112$-vertex counterexample을 제시했고, Jooken이 이어서 human-checkable proof를 공개했다.
\end{solution}

\begin{problem}{5}
다음 two-player zero-sum game에서 row player는 payoff를 최대화하고 column player는 최소화한다. 두 사람은 동시에 한 행과 한 열을 선택하며, row player의 payoff matrix는
\[
    \mathbf{M}
    =
    \begin{pmatrix}
        3&0\\
        1&2
    \end{pmatrix}
\]
이다. Row player가 첫 번째 행을 선택하는 optimal probability $p$, column player가 첫 번째 열을 선택하는 optimal probability $q$, 그리고 game value $v$를 구하시오.
\end{problem}

\begin{solution}
Row player가 첫 번째 행을 probability $p$로 선택한다고 하자. Column player가 첫 번째 열을 고를 때와 두 번째 열을 고를 때의 expected payoff는 각각
\[
    1+2p,
    \qquad
    2-2p
\]
이다. 두 값을 같게 두면
\[
    1+2p=2-2p,
\]
이므로
\[
    p=\frac14,
    \qquad
    v=\frac32.
\]

이제 column player가 첫 번째 열을 probability $q$로 선택한다고 하자. Row player가 첫 번째 행과 두 번째 행을 고를 때의 expected payoff는 각각
\[
    3q,
    \qquad
    2-q
\]
이다. 두 값을 같게 두면
\[
    3q=2-q,
\]
이므로
\[
    q=\frac12.
\]
따라서
\[
    \boxed{
        p=\frac14,
        \qquad
        q=\frac12,
        \qquad
        v=\frac32
    }.
\]
\end{solution}

\begin{problem}{6}
$\alpha,\beta,\gamma$가 다항식
\[
    x^3-3x+1
\]
의 세 근일 때, 다음 값을 구하시오.
\[
    \frac{1}{1-\alpha}
    +\frac{1}{1-\beta}
    +\frac{1}{1-\gamma}.
\]
\end{problem}

\begin{solution}
Vi\`ete's formulas에 의해
\[
    \alpha+\beta+\gamma=0,
    \qquad
    \alpha\beta+\beta\gamma+\gamma\alpha=-3.
\]
또한 $1$은 주어진 다항식의 근이 아니므로 분모는 모두 $0$이 아니다. 통분하면 분자는
\[
    \begin{aligned}
        &(1-\beta)(1-\gamma)
        +(1-\gamma)(1-\alpha)
        +(1-\alpha)(1-\beta)\\
        &\qquad
        =3-2(\alpha+\beta+\gamma)
        +(\alpha\beta+\beta\gamma+\gamma\alpha)
        =0.
    \end{aligned}
\]
따라서
\[
    \boxed{0}.
\]
\end{solution}

% Verified against the 2026 International Congress of Basic Science medal announcements.
\begin{problem}{7}
2026 International Congress of Basic Science에서 arithmetic geometry의 transformative contributions로 Shou-Wu Zhang에게 수여된 mathematics medal의 이름을 쓰시오.
\end{problem}

\begin{solution}
Shou-Wu Zhang이 2026 International Congress of Basic Science에서 받은 mathematics medal은
\[
    \boxed{\text{Andrew Wiles Medal}}
\]
이다. 수상 citation은 arithmetic geometry에 대한 transformative contributions를 강조한다.
\end{solution}

\begin{problem}{8}
대칭군 $S_6$가 $\{1,2,\ldots,6\}$을 두 원소씩 세 쌍으로 분할한 partition들의 집합 위에 자연스럽게 작용한다고 하자. Partition
\[
    \mathcal{P}
    \coloneqq
    \bigl\{\{1,2\},\{3,4\},\{5,6\}\bigr\}
\]
의 stabilizer $(S_6)_{\mathcal{P}}$의 order를 구하시오.
\end{problem}

\begin{solution}
$S_6$의 orbit은 $6$개의 원소를 두 원소씩 세 쌍으로 나누는 모든 partition들로 이루어진다. 그 개수는
\[
    \frac{6!}{2^3\,3!}=15.
\]
Orbit--stabilizer theorem에 의해
\[
    \left|(S_6)_{\mathcal{P}}\right|
    =
    \frac{|S_6|}{|S_6\cdot\mathcal{P}|}
    =
    \frac{6!}{15}
    =48.
\]
따라서
\[
    \boxed{48}.
\]
\end{solution}

% Verified against the Courant Institute announcement of Peter D. Lax's death on May 16, 2025.
\begin{problem}{9}
2025년 5월 16일 $99$세로 별세한 수학자를 쓰시오. 이 수학자는 partial differential equations, conservation laws, numerical analysis에 큰 업적을 남겼으며, Lax pair와 Lax equivalence theorem에 이름을 남겼고 2005 Abel Prize를 수상했다.
\end{problem}

\begin{solution}
주어진 설명에 해당하는 수학자는
\[
    \boxed{\text{Peter D. Lax}}
\]
이다. Courant Institute는 그가 2025년 5월 16일 $99$세로 별세했다고 발표했다.
\end{solution}

\begin{problem}{10}
$z\in\C\setminus\{0\}$가
\[
    z+\frac{1}{z}=1
\]
을 만족한다. 다음 값을 구하시오.
\[
    z^{2026}+\frac{1}{z^{2026}}.
\]
\end{problem}

\begin{solution}
주어진 식은
\[
    z^2-z+1=0
\]
과 동치이므로
\[
    z=e^{\pm i\pi/3}.
\]
따라서
\[
    z^m+z^{-m}
    =
    2\cos\left(\frac{m\pi}{3}\right).
\]
$2026\equiv4\pmod6$이므로
\[
    \boxed{
        z^{2026}+z^{-2026}
        =
        2\cos\left(\frac{4\pi}{3}\right)
        =-1
    }.
\]
\end{solution}
